\section{Infrastructure}
\label{system-architecture-and-implementation:infrastructure}

The benchmarking framework runs within a single Azure resource group, \texttt{doppa}, provisioned in Norway East so that every metered resource shares one billing and administrative boundary. The Azure Databricks workspace that hosts the Apache Sedona clusters is the exception, placed in Sweden Central instead. Norway East lacks sufficient \texttt{Standard\_D4s\_v3} quota for the multi-node clusters required by the national-scale spatial join, so the workspace is created where that quota is available. This split is the only source of cross-region traffic in the system: the Sedona executors in Sweden Central read \acrshort{adls} in Norway East, whereas every other configuration keeps its compute and storage within the same region and tenant. The resulting egress is the single billed transfer in the framework, and its rate enters the cost model in Section~\ref{system-architecture-and-implementation:measurement-instrumentation-and-cost-model:cost-model-implementation}.

\subsection{Compute platforms}
\label{system-architecture-and-implementation:infrastructure:compute-platforms}

Each single-node experiment runs inside its own \acrshort{aci} container group. Every benchmark is a short-lived job that needs an isolated execution unit with no cross-experiment state, and \acrshort{aci} provides this with no orchestrator or node pool persisting between runs. The orchestrator creates one container group per experiment with a restart policy of \texttt{Never}, streams its logs to completion, and deletes the group before the next experiment is dispatched, so each measurement starts from a fresh process on a freshly scheduled host. Query containers are sized at 4 vCPU and \SI{16}{\gibi\byte} of memory, matching the DuckDB and PostGIS client allocations and giving all single-node configurations an identical compute envelope. The outside-cluster orchestrator runs in its own \acrshort{aci} container group sized at 2 vCPU and \SI{2}{\gibi\byte}, since its workload is the serial dispatch loop of Section~\ref{system-architecture-and-implementation:orchestration-and-pairing}, not query execution.

Apache Sedona runs on Azure Databricks classic compute, with one fresh job cluster provisioned per Sedona experiment and terminated in a \texttt{finally} block after the timed iterations complete. Classic clusters expose node type and worker count, both factors of the \acrshort{rq}2 design, whereas serverless compute does not. The driver and every worker run on \texttt{Standard\_D4s\_v3} virtual machines (4 vCore, \SI{16}{\gibi\byte}). Each Sedona node therefore carries the same per-node resource envelope as a single-node \acrshort{aci}, and variation between cluster sizes traces to worker count alone. The runtime is pinned to Databricks Runtime 15.4 \acrshort{lts} (\texttt{15.4.x-scala2.12}). Cluster provisioning, library installation, and the post-run termination all fall outside the timed window, per the warmup discipline established in Section~\ref{research-design-and-methodology:overview-and-experimental-design:replication-and-randomization}.

PostGIS runs against an Azure Database for PostgreSQL Flexible Server provisioned in Norway East as a General Purpose \texttt{Standard\_D4ads\_v5} instance (4 vCores, \SI{16}{\gibi\byte} of memory, AMD processor). Storage is Premium SSD at the P10 performance tier, sized to \SI{128}{\gibi\byte}, with zonal high availability disabled and storage autogrow off. The General Purpose tier was selected over Burstable so that compute capacity is pre-allocated and billed at a flat hourly rate. A Burstable instance would otherwise introduce credit-driven variability in throughput across the long-running national-scale join. The same instance backs every PostGIS workload at every dataset tier, so all PostGIS results share one server configuration and the cost terms developed in Section~\ref{system-architecture-and-implementation:measurement-instrumentation-and-cost-model:cost-model-implementation} apply uniformly.

\subsection{Storage and identity}
\label{system-architecture-and-implementation:infrastructure:storage-and-identity}
% Single \acrshort{adls} account `doppabs`, hierarchical namespace; blob vs dfs endpoints
%   (DuckDB az:// / abfss:// for Sedona). Soft-delete disabled — README IMPORTANT: enabling it
%   makes the Databricks ABFS driver fail with HTTP 409 on dfs.core.windows.net reads. State that
%   as the stated rationale.
% Identity: single user-assigned managed identity `doppa-uami`, Contributor on the RG, AcrPull/
%   AcrPush on the registry; GitHub Actions federated credentials; no secret material baked into
%   images (secrets via --secure-environment-variables from main.py / workflows).
% Container registry `doppaacr`; images pushed by push-containers-to-acr.yml.

\subsection{Determinism and reproducibility}
\label{system-architecture-and-implementation:infrastructure:determinism-and-reproducibility}
% YAML manifest (benchmarks.yml) + per-run seed: main.py shuffles experiments with
%   random.Random(benchmark_run); realizes the randomization of
%   Section~\ref{research-design-and-methodology:overview-and-experimental-design:replication-and-randomization}.
% Run ID = date + 6-char suffix (RUN_ID_LENGTH); benchmark_run indexes repeated suite runs
%   (BENCHMARK_RUNS=30). Bootstrap RNG seeded from (run_id, query_id) via blake2b (monitor_utils
%   _make_bootstrap_rng) — implementation of the reproducible stop point promised in Ch.4.
% Warmup: BENCHMARK_WARMUP_ITERATIONS=5 single-node; Databricks + national-scale joins override
%   warmup_iterations=1 (one warmup is enough on long low-variance runs; more dominates the
%   wall-clock budget — stated CLAUDE.md/README rationale). Cross-ref Ch.4 replication protocol.